\chapter{HPD position reconstruction}
\label{chap:recon}

I need an introduction on the detector geometry before describing how we can make the position reconstruction.

We need to talk about the cluster size distribution and then we can talk about the implementation of different methods for various cluster sizes. We should also first point out that we are always excluding the fourth and successive pixels in an event, because with the given properties of the sensor, the charge collected by the fourth pixel is always below 20 e, while we currently have a higher noise level. This means that we cannot reconstruct the position using this information. 


The main feature of this detector is that the silicon sensor is on top of the readout chip, and its thickness is enough to convert all the energy of a soft X-ray photon inside it, while producing a small diffusion of the produced electrons, with respect to the size of the pixels. This can be seen with a simple calculation, because the transverse diffusion has a sigma of about 40 um/cm**0.5 in silicon, and with a sensor 300 um. This means that for this thickness, the charge cloud has a sigma of about 7 um, smaller than the pixel pitch.

The first method that can be thought to reconstruct the incident position, is to calculate the charge barycenter using the position of the hit pixels and the charge that they collected. It can be shown that this method is not accurate at all, showing a high degree of bias in the reconstruction. This is due to the charge diffusion nature and the electron cloud size. The cloud 98\% of the charge is contained in a 2 sigma circle, which means that a pixel can collect signal only if the photon hits at about this length from the pixel edge, otherwise no charge is collected. Charge barycentering is highly problematic when a photon hits in the edge region, where even if a few electron hit the second pixel, the position would be reconstructed in proximity of the first pixel center. This problem is evident in the case of two.



\section{Centroid}
The cluster centroid is defined as the charge distribution center of gravity, with the assumption that all the charge in a pixel is located in its center. This means that the two position components are given by

\begin{equation}
x = \frac{\sum_i x_i \cdot q_i}{\sum_i q_i}, \quad y = \frac{\sum_i y_i \cdot q_i}{\sum_i q_i}.
\end{equation}

This algorithm doesn't consider how the charge cloud diffuses inside the sensor, and assumes that the charge is linearly divided in the pixels as the position changes. This is clear when considering a two-pixel cluster: if we have 90\% of the charge in a pixel and 10\% in the second one, with this algorithm we get that the photon hit is 0.1 pitch away from the center of the first pixel. In our case, where the charge cloud is smaller than the pixel pitch, it is obvious that this statement is clearly wrong. Fig. \ref{fig:2_pix_algorithms} shows how the reconstructed position varies as a function of the charge in the second pixel.

Centroid's performance improves with larger charge cloud, as in the case of IXPE, where an event cluster is made of tens of pixels. With this detector geometry, it is not a possibility to use it to reconstruct the position of an event.

[Insert some plots of the performance, the radial distance distribution]


\section{Eta function}

An algorithm that can be implemented is the so-called eta function, which consists in studying the relation between the incident position and $\eta$, which is the ratio between the charge collected by a pixel and the total charge. The definition of $\eta$ can vary from case to case, as we will see later, but the concept is always the same: find an analytical model that describes the charge diffusion, calibrate it using a Monte Carlo simulation of the detector, and then use it to reconstruct positions.

To model the charge diffusion, we know that the diffusion is a gaussian process in the plane of the detector. With this information, we can calculate the charge collected by the pixels integrating the two-dimensional Gaussian distribution, centered in the incident position, in the region outlined by each pixel. An analytical solution of this two-dimensional problem is not possible, and in some cases even useless, as in the case of one or two-pixel events, but it is possible to use some approximations and still obtain good results.

In the case of one-pixel events, it is straightforward to conclude that the only possibility is to reconstruct the event in the center of the pixel. In the following we will see how it is possibile to implement this algorithm for two-pixel and three-pixel events.

\subsection{Two-pixel events}

\begin{figure}[htbp]
    \centering
    \begin{tikzpicture}[scale=4]
    \hexpix{\hexsize}
    \node[anchor=south] at (0, 0) {$q_1$};
    \begin{scope}
        \clip (0,0) -- (30:\hexsize) -- (-30:\hexsize) -- cycle;
        \fill[pattern=north east lines, pattern color=gray] (-1,-1) rectangle (1,1);
    \end{scope}
    \hexpix[xshift=1cm]{\hexsize}
    \node[anchor=south] at (1, 0) {$q_2$};
    \draw[|-|] (0, 0.75) -- (1, 0.75);
    \node[anchor=south] at (0.5, 0.75) {$p$};
    \draw[->] (0, 0) -- (1.75, 0);
    \node[anchor=west] at (1.75, 0) {$r$};
    \end{tikzpicture}
    \caption{Illustration of the two-pixel event topology. Assuming that  $q_2 \leq q_1$, 
    the hatched triangle represents the portion of the phase space in terms of the 
    true impact point originating this topology.}
    \label{fig:two_pixel_event}
\end{figure}

For two-pixel events, we can parametrize the problem in terms of the variable $\eta$, defined as

\begin{equation}
\eta = \frac{q_2}{q_1 + q_2}, \quad q_1 \geq q_2
\end{equation}

which describes the fraction of charge collected by the pixel with the lowest charge. With this ordering, $\eta$ is in the range $[0, 0.5]$. With two pixels, the only position component that can be reconstructed lays on the line that connects the pixels' centers. We can call this component $r$.

To find the model that describes this situation, we can consider the special case of the photon hitting exactly on the line connecting the centers. This assumption reduces the problem to a one-dimensional problem that can be easily solved. As said before, we can calculate the fraction of the integral that falls inside the second pixel, and relate it to the incident position $r$. Considering that the charge cloud is smaller than the pixel dimension, we can consider that all the charge is collected by these two pixels, thus we can calculate that

\begin{equation}
\eta(r) = \int ^{\infty}_{p/2} \frac{1}{\sqrt{2\pi}\sigma} \exp \left[ - \frac{(r' - r)^2}{2\sigma^2}  \right] dr' = 1 - \Phi \left( \frac{r - p/2}{\sigma} \right),
\end{equation}

where $\sigma$ is the transeverse diffusion sigma of electrons in the detector sensor and $\Phi(x)$ is the cumulative distribution of the standard normal distribution. Inverting this relation, we can find $r$ as a function of $\eta$

\begin{equation}
r = \frac{p}{2} - \sigma \Phi^{-1}(1 - \eta) = \frac{p}{2} + \sigma \Phi^{-1}(\eta),
\label{eq:eta_2pix}
\end{equation}

where it has been used that $\Phi^{-1}(1 - \eta) = -\Phi^{-1}(\eta)$ for the simmetry of the Gaussian distribution. The relation that has been found is non-linear, and considers exactly the nature of the diffusion geometry, contrary to the centroid, that in case of two pixels gives $r = p \eta$. The two relations are shown in Fig. \ref{fig:centroid_eta_2pix}, with different values of the transverse diffusion sigma. The smaller the diffusion is, the greater the difference between the two models is.

\subsection{Three-pixel events}

An three-pixel event allows to reconstruct the position on a two-dimensional space, contrary to one and two-pixel events. To model the problem, it is first necessary to find the optimal coordinates to parametrize it. One possibility is to define polar coordinates $(r, \theta)$, with $r$ that ranges between $[0, \frac{p}{\sqrt{3}}]$ and $\theta$ between $[-\frac{\pi}{6}, \frac{\pi}{6}]$, with $\theta = 0$ passing in the midpoint between the outer pixels, as shown in Fig. \ref{fig:three_pixel_event}. As for two-pixel events, pixels' charge are in the order $q_1 \geq q_2 \geq q_3$, and we define outer pixels the two with the least charge. 

\begin{figure}[htbp]
    \centering
    \begin{tikzpicture}[scale=3]
    \hexpix{\hexsize}
    \node[anchor=south] at (0, 0) {$q_1$};
    \draw[|-|] (0, 0.75) -- (1, 0.75);
    \node[anchor=south] at (0.5, 0.75) {$p$};
    \begin{scope}
        \clip (0,0) -- (0:\hexsize) -- (-30:\hexsize) -- cycle;
        \fill[pattern=north east lines, pattern color=gray] (-1,-1) rectangle (1,1);
    \end{scope}
    \hexpix[xshift=1cm]{\hexsize}
    \node[anchor=south] at (1, 0) {$q_2$};
    \hexpix[xshift=0.5cm, yshift=-0.866cm]{\hexsize}
    \node[anchor=north] at (0.5, -0.866) {$q_3$};
    \draw[dashed] (0, 0) -- (1, 0);
    \draw[->] (0, 0) -- (1.5, -0.866);
    \draw[dashed] (0, 0) -- (0.5, -0.866);
    \draw[dashed] (0.5, -0.866) -- (1, 0);
    \draw[dashed] (0, 0) -- (1.7, -0.45);
    \node[anchor=west] at (1.5, -0.866) {$r$};
    \draw[->] (1.45, -0.827) arc[start angle=-30, end angle=-10, radius=1.2];
    \node[anchor=west] at (1.55, -0.65) {$\theta$};
    \end{tikzpicture}
    \caption{Illustration of the three-pixel event topology. Assuming that the ordering
    of the charges is $q_3 \leq q_2 \leq q_1$, the hatched triangle represents again the 
    portion of the phase space in terms of the true impact point originating this topology.}
    \label{fig:three_pixel_event}
\end{figure}

It is now necessary to find the charge coordinates associated to the spatial polar coordinates. For the radial component, we can think of a simplified problem, with the photon hitting at an angle of $0^{\circ}$ with respect to the axis passing through the first pixel vertex. In this configuration, we can think the two outer pixels as a single pixel, whose total charge is the sum of their charge. Thus this problem is reduced to a one-dimensional problem similar to the two-pixel problem, but with a new charge coordinate, defined as

\begin{equation}
\eta_+ = \eta_2 + \eta_3.
\end{equation}

To complete the model, it is first necessary to consider that this time the integral is carried from the vertex of the hexagon, which is located at $r=\frac{p}{\sqrt{3}}$, hence the offset of the model is shifted in this position.

Another thing that has to be tweaked is the normalisation of $\eta_+$. If a photon strikes exactly on the vertex of the hexagon, the charge is equally distributed between the three pixels, thus $\eta_+ = 2/3$. According to Eq. \ref{eq:eta_2pix} with the new offset, this value of $\eta_+$ is not reconstructed on the vertex, but in a point outside the first hexagon. This can be easily corrected with a normalisation factor to impose this constraint. The final relation between $r$ and $\eta_+$ is

\begin{equation}
r = \frac{p}{\sqrt{3}} + \sigma \Phi^{-1} \left( \frac{3}{4}\eta_+ \right).
\label{eq:eta_3pix}
\end{equation}

To model the angular component, we can consider what happens in the approximation of small angles. If the angle increases a bit, the charge in the pixel with the most charge remains the same, while the charge in the second pixel increases and that in the third pixel decreases of the same amount. This suggests that the ideal charge coordinate to describe the problem is the charge difference of the two outer pixels, or a combination of that. We can call this coordinate $\eta_-$, and is defined as

\begin{equation}
\eta_- = \frac{\eta_2 - \eta_3}{\eta_2 + \eta_3},
\end{equation}

which is the charge difference of the outer pixels normalized to their total charge. The idea behind this model is similar to the previous cases but there are some differences. We can think of the interesection of the pixels as a plane divided into three orthogonal parts, as shown in Fig. \ref{fig:angular_3pix}. It can be seen that $\eta_-$ depends on the transverse coordinate $y$, specifically it can be found as the difference between the integral of the distribution in the quadrant of the second pixel and the third pixel

\begin{equation}
\eta_- = \frac{\eta_2}{\eta_2 + \eta_3} - \frac{\eta_3}{\eta_2 + \eta_3} = \left[1 - \Phi \left( - \frac{y}{\sigma} \right ) \right] - \Phi \left( - \frac{y}{\sigma} \right),
\end{equation}

which can be inverted to find $y$. Given that we are considering small angles, the transverse coordinate $y$ is related to the angle by the relation $y=r\theta$, thus we have

\begin{equation}
\theta = - \frac{\sigma}{r} \Phi^{-1} \left( \frac{1 - \eta_-}{2} \right) = \frac{\sigma}{r} \Phi^{-1} \left( \frac{1 + \eta_-}{2} \right).
\end{equation}


\begin{figure}[t]
    \centering
    \begin{tikzpicture}[scale=1.5]
    
    \draw[thick] (0,-2) -- (0,2);
    \draw[thick, ->] (0,0) -- (2,0) node[anchor=west] {$r$}; 
    \draw[thick, dashed] (-2,0) -- (0,0);
    
    \draw[thick] (-0.3,0.3) circle (0.7);
    \draw[thick, dashed, <->] (-0.3, 0.0) -- (-0.3, 0.3) node[anchor=east] {$y$};
    \fill (-0.3,0.3) circle (0.03); 

    \node at (-1.5,1) {Pixel 1};
    \node at (1,1) {Pixel 2};
    \node at (1,-1) {Pixel 3};

    \end{tikzpicture}
    \caption{Illustration of the approximation made for the calculation of the  }
    \label{fig:angular_3pix}
\end{figure}

\subsection{Calibration}

To calibrate the models with the actual response of the detector, it is required to use Monte Carlo simulations to find the best-fit paramenters of the models. All of these models have been obtained with some approximations, and even if $\sigma$ has the same meaning for all of them, real data does not always respect the geometrical assumptions that have been made, leading to differences between the expected value of the parameters and their best-fit value.

The calibration procedure is performes with a generated Monte Carlo dataset with no electronic noise. Between all the events, only those with two or three pixels above a given threshold are considered. The threshold is chosen based on the expected electronic noise of the detector, and for this analysis it has been set to twice the noise sigma. The threshold is necessary in the reconstruction phase, as it is necessary to cut pixels that have a high probablity of only having noise charges, and for two-pixel events it reflects in a minimum value of the $\eta$ coordinate.

To perform the calibration, each component of the two event classes are binned into a two-dimensional space, with the charge coordinate on the x-axis and the spatial component on the y-axis. Then, a profiling is made calculating the median of the position distribution for each value of the charge component. The choice to consider the median instead of the mean was made to limit the influence of the distribution tails on the position estimator. A least-squares fit is performed on the data using the respective models.

[Show some examples]



\section{Maximum likelihood estimator}
This method is the generalization of the eta function algorithm, and uses the charge diffusion model and  the pixel readout probability density function. Given a photon of energy $E$ that strikes a pixel in the coordinate $(x_0, y_0)$, at a first approximation the \textit{pdf} of the charge in the $i$-th pixel is a Gaussian distribution

\begin{equation}
p(q_i; f_i E , \sigma^2_n) = \frac{1}{\sqrt{2\pi}\sigma_n} \exp \left[- \frac{(q_i - f_i E)^2}{2\sigma_n^2} \right],
\label{eq:pdf_gauss}
\end{equation}

where $f_i$ represents the fraction of the total charge deposited in the $i$-th pixel, and depends on the photon impact position. This quantity could in principle be calculated integrating the two-dimensional Gaussian distribution centered in $(x_0, y_0)$ on each of the seven pixels of the cluster

\begin{equation}
f_i (x_0, y_0) = \iint_{P_i} \frac{1}{2\pi \sigma^2} \exp \left[ - \frac{(x - x_0)^2 + (y - y_0)^2}{2\sigma^2} \right] \,dx\,dy,
\end{equation}

where $P_i$ denotes the region of the $i$-th pixel. In practice, these calculations are not easy to do and use in an analysis software, but we will go into details of this part later. Assuming we know what is the value of $f_i$ for each pixel given an incident position $(x_0, y_0)$, we can write the likelihood of  our measurement. The quantities that we measured are the charges collected by the seven pixels, thus the likelihood is

\begin{equation}
\mathcal{L}(x, y, E) = \prod_{i=1}^7 p(q_i; f_iE, \sigma_n^2).
\end{equation}

Here we are assuming to know the noise distribution, which is a true, as it will be described in Ch. \ref{chap:noise}. At this point, we can write the negative log-likelihood ignoring the constant terms and minimize it to find the best-fit values

\begin{equation}
- \log \mathcal{L} = \sum_{i=1}^7 \frac{(q_i - f_iE)^2}{\sigma_n^2}.
\end{equation}

As said at the beginning of this section, this result has been obtained with the approximation that the charge \textit{pdf} is a Gaussian distribution. This is true when the counts of the pixel are above zero, but the readout chip clips values below zero counts. This clipping modifies Eq. \ref{eq:pdf_gauss}, and now we have to take into account that if $q_i$ is zero, we need to calculate the probability of having a value smaller than zero, meaning we have to integrate the Gaussian centered in $f_i E$ up to zero. Now we have a mixed probability distribution

\begin{equation}
p(q_i; f_i E , \sigma^2_n) = 
\begin{cases}
\frac{1}{\sqrt{2\pi}\sigma_n} \exp \left[- \frac{(q_i - f_i E)^2}{2\sigma_n^2}\right] & q_i > 0\\
\Phi \left(- \frac{f_iE}{\sigma_n} \right) & q_i = 0
\end{cases}
\end{equation}

This distribution now takes into account all the physical and readout aspects of the chip, and the resulting negative log-likelihood can be minimized with any scientific library to find the impact position and energy best-fit values.

An aspect that was left behind is how we calculate $f_i$ for each position. As said before, it could be calculated integrating the two-dimensional Gaussian centered in $(x, y)$ on all of the pixels, but it's not computationally practical. A different way to do this task is to use Monte Carlo simulations generated with absent electronic noise, in order to take into account only how the charge cloud spread across the pixels. The simulation returns the exact amount of charge for each of the seven pixels in the cluster, and dividing for the total charge we can calculate the charge fraction. Simulating a dataset with photon impact positions spread uniformly across the central pixel, it is possible to sample $f_i(x, y)$ for each pixel by creating a two-dimensional grid over the hexagon and calculating the average between all the events in each bin. The bin width has to be of the same magnitude of the expected spatial resolution, in order to be robust to the noise in the real data.
[Plot delle griglie]


























